\section{The cleartext oracle, and how many power iterations are needed}
\label{sec:tech:oracle}

Every private component in this project is checked against a cleartext twin.
The factorisation oracle is therefore the first thing built, before any
cryptography, and it doubles as baseline B1, the speed-of-light reference
against which recommendation quality is measured.

\subsection{Construction}
\label{sec:tech:oracle:construction}

The oracle implements $\ApproxFactor$ as \nudgeplain{} states it, rather than
calling a singular value decomposition. This is deliberate. The private
implementation in the next phase must reproduce this procedure step by step
under secret sharing, so the reference has to match the \emph{structure} of
the protocol and not merely its output. Concretely, for each of $\dime$
components a random vector is repeatedly multiplied by
$\ratings^{\top}\ratings$, made orthogonal to the components already found,
and normalised. Each converged component is then revealed in the clear and
becomes a row of $\itemembed$.

Revealing each row before computing the next is what keeps the
orthogonalisation cheap, because it is then a Gram--Schmidt step against
public vectors. It is also the reason the item embedding matrix is public to
every server, which is the premise of the leakage analysis in this project.

A singular value decomposition is still computed, but as a \emph{cross-check}
rather than as the implementation. Power iteration on
$\ratings^{\top}\ratings$ converges to the top-$\dime$ right singular
subspace, so the largest principal angle between the two subspaces should
approach zero. That check costs two lines and it caught a genuine question on
the first day.

\subsection{Quality is flat in the iteration count}
\label{sec:tech:oracle:ell}

At the default of $\iters = 10$ inner rounds the subspace angle is
$9.66 \times 10^{-1}$ radians, which is nowhere near zero. Sweeping $\iters$
resolves what that means.

\begin{table}[h]
\centering
\small
\caption{MovieLens-100K, split \texttt{u1}, $\dime = 16$. The subspace angle
falls by five orders of magnitude while recommendation quality does not move.
Raw data in \texttt{bench/results/oracle\_ell\_sweep.jsonl}.}
\label{tab:tech:ell}
\begin{tabular}{@{}rrrr@{}}
\toprule
$\iters$ & subspace angle (rad) & nDCG@20 & Recall@20 \\
\midrule
10  & $9.659 \times 10^{-1}$ & 0.4536 & 0.4693 \\
25  & $4.170 \times 10^{-1}$ & 0.4545 & 0.4677 \\
50  & $1.834 \times 10^{-1}$ & 0.4525 & 0.4650 \\
100 & $6.966 \times 10^{-2}$ & 0.4523 & 0.4662 \\
200 & $3.604 \times 10^{-3}$ & 0.4526 & 0.4664 \\
400 & $6.652 \times 10^{-6}$ & 0.4526 & 0.4664 \\
\bottomrule
\end{tabular}
\end{table}

\begin{figure}[h]
\centering
\includegraphics[width=0.68\linewidth]{figures/fig_ell}
\caption{The same data as Table~\ref{tab:tech:ell}, plotted. Note the axes:
nDCG@20 spans a range of 0.002 while the subspace angle spans five orders of
magnitude. Since $\iters$ multiplies the only interactive cost in private
training, the flat curve is what makes a fortyfold communication saving free.
Regenerated by \texttt{mingw32-make figures}.}
\label{fig:ell}
\end{figure}

Two conclusions follow, and they are separate.

The first is that the implementation is correct. The angle falls
monotonically to $6.7 \times 10^{-6}$ radians, so the procedure does converge
to the right subspace. The large angle at $\iters = 10$ is slow convergence
rather than a defect. The spectrum explains the rate: the ratio of the
sixteenth to the seventeenth squared singular value is $1.032$, a gap of three
percent, and power iteration converges as a power of that ratio. The first
component is easy, with a ratio of $6.44$. The sixteenth is not.

This distinction was nearly lost. A failing cross-check invites the reaction
of removing the cross-check, which would have converted a correct
implementation into a silently wrong one.

The second conclusion is the useful one. Across a fortyfold increase in
$\iters$, nDCG@20 stays between $0.4523$ and $0.4545$, a spread of about half
a percent with no monotone trend. Ranking does not require the singular
vectors, only a subspace good enough that the ordering of scores is stable.

\subsection{Why this matters for private training}
\label{sec:tech:oracle:consequence}

Under two-out-of-three replicated sharing the matrix--vector products are
non-interactive and therefore free. Truncation and normalisation are the
entire communication cost of training, and they run once per inner iteration,
so communication is linear in $\iters$.

The measurement therefore says that private training can run at $\iters = 10$
and the fortyfold saving in communication costs nothing in recommendation
quality. That is a design decision resting on evidence rather than on the
default value in the design draft, and it also discharges an instruction the
draft had already given and which had not been acted on, namely to measure
the residual against the cleartext oracle and report iterations against
quality.

The caveats are stated plainly. This is one dataset, one split, one seed and
one embedding dimension. The flatness may not survive at MovieLens-1M or at
$\dime = 32$, and it will be re-measured before the final report relies on it.
The metric uses binary relevance at a rating of four or above, so these
numbers are a within-project baseline and are not directly comparable to the
figures \nudgeplain{} reports on Netflix data under its own convention.
